% Quelle: 5. Zyklische Codes.pdf
% Art: Vorlesung
% Themen: zyklische Codes, Wörter als Polynome, Multiplikation in F[X]/(X^n-1), Ideal, Generatorpolynom, Generatormatrix, Kodieren, Syndrom, n-te Einheitswurzeln, Kreisteilungsklassen, Zerlegung von X^n-1, Kontrollpolynom, Kontrollmatrix, reziprokes Polynom, binärer Golay Code G23, ternärer Golay Code G11, Hamming-Schranke, perfekte Codes

\section{5. Zyklische Codes}

\subsection{5.1 Definition und erste Beispiele}

\begin{definition}
Ein linearer Code $C \subseteq F^n$ heißt \textbf{zyklisch}, wenn:
\[
c_0\ c_1\ \dots\ c_{n-2}\ c_{n-1} \in C
\quad\Rightarrow\quad
c_{n-1}\ c_0\ \dots\ c_{n-3}\ c_{n-2} \in C
\]
d.h.\ eine zyklische Verschiebung eines Codewortes ist wieder ein Codewort.
\end{definition}

\begin{beispiel}[Beispiele]\leavevmode
\begin{enumerate}[label=\alph*)]
\item Der binäre Code $\{000, 101, 011, 110\}$ ist zyklisch.
\item Der Code
\[
\{0000, 1001, 0110, 1111\} \subseteq \{0,1\}^4
\]
ist nicht zyklisch. Vertauschen der 3.\ und 4.\ Koordinate liefert einen zyklischen Code.
\item $C_3$ = binärer Paritätscode mit 3 Informationsbits und 1 Kontrollbit
\[
C_3 = \{\, 0000, 1100, 0110, 0011, 1001, 0101, 1010, 1111 \,\} \subseteq \{0,1\}^4
\]
ist zyklisch.
\item Das Codewort $\mathbf{g} = 1101000$ kann man durch zyklische Verschiebung zu einer Generatormatrix eines $[7,4,3]_2$ Codes ergänzen
\[
G = \begin{pmatrix}
1 & 1 &   & 1 &   &   &   \\
  & 1 & 1 &   & 1 &   &   \\
  &   & 1 & 1 &   & 1 &   \\
  &   &   & 1 & 1 &   & 1
\end{pmatrix}.
\]
\begin{itemize}
\item[-] Welcher uns bekannte Code ist das?
\item[-] Wie findet man eine Kontrollmatrix?
\item[-] Welche weitere interessante zyklische Codes gibt es?
\item[-] Die 1.\ Zeile enthält bereits sämtliche Informationen zu diesem Code. Welche Mathematik kann man nutzen?
\end{itemize}
\end{enumerate}
\end{beispiel}

\subsection{5.2 Wörter als Polynome}

\[
\mathbf{a} = a_0\, a_1 \dots a_{n-1} \quad\to\quad a(X) = a_0 + a_1 X + \dots + a_{n-1} X^{n-1}
\]

\begin{enumerate}
\item Isomorphismus \qquad $F^n \to F_n[X]$
\end{enumerate}
Wünschenswert ist
\begin{enumerate}[start=2]
\item $X * a(X) = a_{n-1} + a_0 X + a_1 X^2 + \dots + a_{n-2} X^{n-1}$
\end{enumerate}
Die Multiplikation \glqq$*$\grqq{} erfolgt aber nicht in $F[X]$, sondern in $F[X]/(X^n - 1)$.

\begin{bemerkung}[Abkürzung]
\[
F_n[X] := F[X]/(X^n - 1)
\]
\begin{tabular}{lll}
\glqq$\cdot$\grqq & Multiplikation in $F[X]$ & z.B.\ $1 \cdot X^n = X^n$, \quad $X \cdot X^n = X^{n+1}$ \\
\glqq$*$\grqq & Multiplikation in $F_n[X]$ & $1 * X^n = 1$, \quad $X * X^n = X$, \quad usw.
\end{tabular}
\end{bemerkung}

\paragraph{Division mit Rest}
\begin{align*}
& a(X) = q(X)\cdot(X^n - 1) + r(X), \quad \operatorname{grad} r(X) < n \\
\Rightarrow\quad & a(X) \equiv r(X) \quad (\bmod\ X^n - 1) \\
\Rightarrow\quad & a(X) = r(X) \quad \text{in} \quad F_n[X]
\end{align*}

\paragraph{Multiplikation in $F^n$}
\begin{align*}
010\dots0 * a_0\, a_1 \dots a_{n-1} &\cong X * (a_0 + a_1 X + \dots + a_{n-1} X^{n-1}) \\
&= a_0 X + a_1 X^2 + \dots + a_{n-2} X^{n-1} + a_{n-1} X^n \\
&= a_0 X + a_1 X^2 + \dots + a_{n-2} X^{n-1} + a_{n-1}
\end{align*}
\[
\Rightarrow\quad 010\dots0 * a_0\, a_1 \dots a_{n-1} = a_{n-1}\, a_0\, a_1 \dots a_{n-2}
\]
Die Buchstaben von $a$ werden zyklisch eins nach rechts verschoben.

\paragraph{Multiplikation $\mathbf{a} * \mathbf{b}$}
\begin{align*}
& \mathbf{a} * \mathbf{b} \cong a(X)\cdot b(X) = q(X)\cdot(X^n - 1) + r(X) \equiv r(X) \quad (\bmod\ X^n - 1) \\
& \mathbf{a} * \mathbf{b} = \mathbf{r}
\end{align*}

\begin{beispiel}[$F = F_5$, $n = 4$]
\begin{align*}
0012 * 2314 &= (X^2 + 2X^3) * (2 + 3X + X^2 + 4X^3) \\
&= X^2 * (2 + 3X + X^2 + 4X^3) + X^3 * (4 + X + 2X^2 + 3X^3) \\
&= 2102
\end{align*}
\end{beispiel}

Sei $\mathbf{v} \in F^n$. $\{\lambda \mathbf{v} \mid \lambda \in F\} \subseteq F^n$ ist ein Unterraum (eine Gerade durch Null).

Von $\mathbf{g}$ erzeugtes \textbf{Ideal}\quad $\{\mathbf{a} * \mathbf{g} \mid \mathbf{a} \in F^k\} \subseteq F^n$ \qquad ($\mathbf{g} \in C$ fest, $k = \operatorname{Dim} C$)

\begin{satz}
Für einen zyklischen Code $C$ gilt
\[
\mathbf{g} \in C \;\Rightarrow\; \mathbf{a} * \mathbf{g} \in C \quad \text{für alle } \mathbf{a} \in F^k
\]
d.h.
\[
F^k * \mathbf{g} \subseteq C.
\]
\end{satz}

\begin{proof}
\[
\mathbf{a} = a_0\, 0\dots0 + 0\, a_1\, 0\dots0 + \dots + 0\dots0 a_{k-1}
\]
also da $C$ linear ist
\begin{align*}
\mathbf{a} * \mathbf{g} &= a_0\, a_1 \dots a_{k-1} * \mathbf{g} \\
&= a_0\, 0\dots0 * \mathbf{g} + 0\, a_1\, 0\dots0 * \mathbf{g} + \dots + 0\dots0 a_{k-1} * \mathbf{g} \\
&\in C \qedhere
\end{align*}
\end{proof}

\subsection{5.3 Generatorpolynom $g(X)$ und Generatormatrix $G$}

Sei $G$ eine Generatormatrix von $C$, $k = \operatorname{Dim} C$, also
\[
C = \mathbf{a}\, G = \{\, a_0 \mathbf{g}_0 + \dots + a_{k-1} \mathbf{g}_{k-1} \mid (a_0 \dots a_{k-1}) \in F^k \,\}
\]
Gibt es ein Codewort $\mathbf{g} \in C$, so dass $F^k * \mathbf{g} = C$?\\
Wenn JA: Ersetzen die Generatormatrix $G$ durch das Wort $\mathbf{g} \in C$ bzw.\ das Generatorpolynom $g(X)$.

\begin{beispiel}[$C_3$, $k = \operatorname{Dim} C_3 = 3$, $F^k = \{000, 100, 010, 001, 111, 110, 011, 101\}$]\leavevmode
\begin{enumerate}[label=\alph*)]
\item $\mathbf{g} = 0110$
\begin{align*}
\Rightarrow\quad C_3 &= \{\, 0000,\ 0110,\ 0011,\ 1001,\ 1100,\ 0101,\ 1010,\ 1111 \,\} \\
&= \{\, 000 * \mathbf{g},\ 100 * \mathbf{g},\ 010 * \mathbf{g},\ 001 * \mathbf{g},\ 111 * \mathbf{g},\ 110 * \mathbf{g},\ 011 * \mathbf{g},\ 101 * \mathbf{g} \,\} \\
&= F^k * \mathbf{g}
\end{align*}
\item $\mathbf{c} = 1010$
\[
\Rightarrow\quad F^k * \mathbf{c} = \{\, 0000, 0101, 1010, 1111 \,\} \subseteq C_3
\]
\item $\mathbf{g} = 1100$
\[
\Rightarrow\quad F^k * \mathbf{g} = C_3 \qquad \text{Überprüfen!!}
\]
\end{enumerate}
\end{beispiel}

\begin{satz}
Sei $\{\mathbf{o}\} \neq C \subseteq F_n[X]$ ein zyklischer Code. Dann gilt
\begin{enumerate}[label=(\arabic*)]
\item Es existiert ein normiertes Polynom $g(X) \neq o(X)$ kleinsten Grades $r$ in $C$.
\[
g(X) = g_0 + g_1 X + \dots + g_{r-1} X^{r-1} + X^r
\]
$g(X)$ ist eindeutig bestimmt.
\item \[ F^k * \mathbf{g} = C \]
\item $g(X)$ ist ein Teiler von $X^n - 1$.
\end{enumerate}
\end{satz}

\begin{proof}\leavevmode
\begin{enumerate}[label=(\arabic*)]
\item Seien $g(X), h(X) \in C$ zwei solche Polynome, $g_r = h_r = 1$.
\begin{align*}
\Rightarrow\quad & g(X) - h(X) \in C, \quad \operatorname{grad}(g - h)(X) < r \\
\Rightarrow\quad & g(X) - h(X) = o(X) \\
\Rightarrow\quad & g(X) = h(X)
\end{align*}
\item Sei $c(X) \in C$, $c(X) \neq o(X)$. Dividieren $c(X)$ mit Rest durch $g(X)$.
\begin{align*}
& c(X) = q(X)\cdot g(X) + r(X) \quad \text{mit } \operatorname{grad} r(X) < r \\
\Rightarrow\quad & r(X) \in C, \qquad \text{da } C \text{ linear} \\
\Rightarrow\quad & r(X) = o(X)
\end{align*}
\item Das ist auch richtig für $c(X) = X^n - 1 = o(X) \in C$.
\[
\Rightarrow\quad X^n - 1 = q(X)\cdot g(X) \qedhere
\]
\end{enumerate}
\end{proof}

\begin{definition}
Das eindeutig bestimmte Polynom
\[
g(X) = g_0 + g_1 X + \dots + g_{r-1} X^{r-1} + X^r
\]
mit den Eigenschaften (1), (2), (3) heißt \textbf{Generatorpolynom} von $C \subseteq F^n$.
\end{definition}

Durchlaufen der Teiler von $X^n - 1$ liefert sämtliche zyklischen Codes $C \subseteq F^n$.

\begin{beispiel}
Bestimmen die zyklischen Codes der Länge $n = 4$ in $\mathbf{F}_2$: \quad $X^4 - 1 = (X + 1)^4$
\begin{align*}
C_0 &= F_0[X] * (1 + X)^4 = F^0 * 0000 = \{\, 0000 \,\} \\
C_1 &= F_1[X] * (1 + X)^3 = F^1 * 1111 = \{\, 0000,\ 1111 \,\} \\
C_2 &= F_2[X] * (1 + X)^2 = F^2 * 1010 = \{\, 0000,\ 1010,\ 0101,\ 1111 \,\} \\
C_3 &= F_3[X] * (1 + X) \phantom{{}^2} = F^3 * 1100 = \{\, \dots \,\} \\
C_4 &= F_4[X] * 1 \phantom{(1 + X)^2} = F^4 * 1000 = F^4
\end{align*}
\end{beispiel}

\begin{satz}
Sei $g(X) = g_0 + g_1 X + \dots + g_{r-1} X^{r-1} + X^r$ das \textbf{Generatorpolynom} von $C$.

Dimension von $C$ \qquad $k + r = n$

Generatormatrix
\[
G = \begin{pmatrix}
g_0 & g_1 & \cdots & g_{r-1} & 1 & & & 0 \\
0 & g_0 & g_1 & \cdots & g_{r-1} & 1 & & \\
 & \ddots & \ddots & \ddots & & \ddots & \ddots & \\
0 & & 0 & g_0 & g_1 & \cdots & g_{r-1} & 1
\end{pmatrix}
=
\begin{pmatrix}
\mathbf{g} \\ X * \mathbf{g} \\ \vdots \\ X^{k-1} * \mathbf{g}
\end{pmatrix}
\]
\end{satz}

\begin{beispiel}[Fortsetzung (Paritätscode $C_3$)]
$g(X) = 1 + X$ ist Generatorpolynom von $C_3$.\\
Tatsächlich ist $k = 3$, $r = 1$, $k + r = 3 + 1 = 4 = n$\\
Die Generatormatrix für obigen Code $C_3$ ist
\[
G_3 = \begin{pmatrix}
1 & 1 & 0 & 0 \\
0 & 1 & 1 & 0 \\
0 & 0 & 1 & 1
\end{pmatrix}
\]
\end{beispiel}

\begin{algorithmus}[Kodieren, $\mathbf{a}$ Klartextwort]\leavevmode
\begin{itemize}
\item[-] Mit Generatormatrix $G$\\
Sei $\mathbf{a} \in F^k$ \quad $\Rightarrow$ \quad $\mathbf{c} = \mathbf{a}\, G$
\item[-] Mit Generatorpolynom $g(X)$\\
Sei $a(X) \in F_k[X]$, $\operatorname{grad} a(X) < k$ \quad $\Rightarrow$ \quad $c(X) = a(X)\cdot g(X)$
\end{itemize}
\end{algorithmus}

\subsection{5.4 Das Syndrom $s(X)$ eines zyklischen Codes}

\begin{definition}
Sei $y(X) \in F^n(X)$ und $g(X)$, $\operatorname{grad} g(X) = r$, das Generatorpolynom von $C$.
Dividieren $y(X)$ mit Rest durch $g(X)$
\[
y(X) = q(X)\cdot g(X) + s(X) \quad \text{mit } \operatorname{grad} s(X) < r.
\]
$s(X)$ heißt \textbf{Syndrom} von $y(X)$.
\end{definition}
In Abschnitt 4.3 über Hamming-Codes haben wir das Syndrom benutzt um zu prüfen, ob $y(X)$ ein Codewort ist und wenn nicht, das nächstgelegene Codewort zu bestimmen.

\subsection{5.5 $n$-te Einheitswurzeln und Kreisteilungsklassen}

\paragraph{Anwendung auf $X^n - 1 \in \mathbf{F}_p[X]$}

Sei $\{\, x \mid x^n - 1 = 0 \,\}$ ist die Menge der $n$-ten Einheitswurzeln von $X^n - 1$.

\begin{satz}
Falls $n$, $p$ teilerfremd, d.h.\ $\ggT(n, p) = 1$, dann gilt
\[
X^n - 1 \text{ besitzt } n \text{ verschiedene Nullstellen.}
\]
\end{satz}

\begin{definition}[Bahn, Kreisteilungsklassen]
\textbf{Bahn} in $\Zm{n}$ \qquad $B_s = \{\, s,\ s\,p,\ s\,p^2,\ \dots,\ s\,p^{d-1} \,\}$, \quad $s\,p^d = s \quad (\bmod\ n)$

$\Zm{n}$ zerfällt in Bahnen, die \textbf{Kreisteilungsklassen.}
\[
\Zm{n} = B_0 \cup B_s \cup B_t \cup \dots
\]
\end{definition}

\begin{satz}\leavevmode
\begin{enumerate}[label=\alph*)]
\item Zu jeder Kreisteilungsklasse $B_s$ existiert ein normiertes, irreduzibles Polynom
\[
m_s(X), \quad d = \operatorname{grad} m_s(X), \qquad X^n - 1 = \prod m_s(X)
\]
\item Falls $n = q \neq p$ eine Primzahl ist, dann haben alle Kreisteilungsklassen $B_s \neq B_0$ dieselbe Länge $d$. Ihre Anzahl ist $\dfrac{q-1}{d}$.
\item Im Fall von $q \equiv \pm 1 \pmod 8$ ist ihre Anzahl 2.
\end{enumerate}
\end{satz}

Anzahl irreduziblen Faktoren von $X^n - 1$ = Anzahl Kreisteilungsklassen.

\begin{beispiel}[$X^n - 1 \in F_2[X]$, $p = 2$]\leavevmode

$n = 5$
\begin{align*}
\Zm{5} &= \{0\} \cup \{1, 2, 4, 3\} \\
X^5 - 1 &= (X + 1)(X^4 + X^3 + X^2 + X + 1)
\end{align*}

$n = 9$
\begin{align*}
\Zm{9} &= \{0\} \cup \{1, 2, 4, -1, -2, -4\} \cup \{3, -3\} \\
X^9 - 1 &= (X^3)^3 - 1 = Y^3 - 1 \quad \text{mit } Y = X^3 \\
&= (Y + 1)(Y^2 + Y + 1) \\
&= (X + 1)(X^2 + X + 1)(X^6 + X^3 + 1)
\end{align*}

$n = 15$
\[
\Zm{15} = \{0\} \cup \{1, 2, 4, -7\} \cup \{3, 6, -3, -6\} \cup \{\pm 5\} \cup \{7, -1, -2, -4\}
\]
irreduzible Polynome in $F_2[X]$ vom Grad 2
\[
X^2 + X + 1
\]
irreduzible Polynome in $F_2[X]$ vom Grad 4
\[
X^4 + X + 1, \quad X^4 + X^3 + 1, \quad X^4 + X^3 + X^4 + X + 1
\]
% UNSICHER: Im Original steht "X^4 + X^3 + X^4 + X + 1" (zweimal X^4); gemeint ist vermutlich X^4 + X^3 + X^2 + X + 1. 1:1 übernommen.
\[
\Rightarrow\quad X^{15} - 1 = (X + 1)(X^2 + X + 1)(X^4 + X + 1)(X^4 + X^3 + 1)(X^4 + X^3 + X^4 + X + 1)
\]
\end{beispiel}

\subsection{5.5* Kontrollpolynom $h(X)$ und Kontrollmatrix $H$}

Sei $g(X)$ das Generatorpolynom von $C$. Dann existiert das Polynom $h(X)$ mit
\begin{align*}
g(X)\cdot h(X) &= X^n - 1 & &\Rightarrow\quad h_k = 1 \\
\operatorname{grad} g(X) + \operatorname{grad} h(X) &= n & &\Rightarrow\quad \operatorname{grad} h(X) = k
\end{align*}

Sei $c(X) \in F_n[X] * g(X) = C$ ein Codewort, also $c(X) = f(X)\cdot g(X)$.
\[
\Rightarrow\quad c(X)\cdot h(X) = f(X)\cdot g(X)\cdot h(X) = f(X)\cdot(X^n - 1) \equiv 0 \quad (\bmod\ X^n - 1)
\]

\begin{definition}
Dieses zu $g(X)$ komplementäre Polynom $h(X)$ heißt \textbf{Kontrollpolynom} von $C$.
\end{definition}

\begin{satz}
Für das Kontrollpolynom
\[
h(X) = h_0 + h_1 X + \dots + h_{k-1} X^{k-1} + X^k
\]
des zyklischen Codes $C$, $k = \operatorname{Dim} C$, gilt
\begin{gather*}
c(X) \in C \quad\Leftrightarrow\quad c(X)\cdot h(X) \equiv 0 \quad (\bmod\ X^n - 1) \\
\operatorname{grad} g(X) + \operatorname{grad} h(X) = r + k = n \qquad \square
\end{gather*}
\end{satz}

\begin{beispiel}[$\mathbf{F}_2$: \ $X^4 - 1 = (X + 1)^4$]
\begin{tabular}{lll}
$C_1 = \{\, 0000,\ 1111 \,\}$, & $g(X) = (1 + X)^3$, & $h(X) = 1 + X$ \\
$C_2 = \{\, 0000,\ 1010,\ 0101,\ 1111 \,\}$, & $g(X) = (1 + X)^2$, & $h(X) = (1 + X)^2$ \\
$C_3 = \{\, 0000,\ 1100, \dots \,\}$, & $g(X) = 1 + X$, & $h(X) = (1 + X)^3$
\end{tabular}
\end{beispiel}

\paragraph{Kontrollmatrix $H$}

\textbf{Ziel}: Finden der $r = n - k$ Zeilen von $H$
\begin{align*}
0 &\equiv c(X)\cdot h(X) \\
&= (c_0 + c_1 X + \dots + c_{n-1} X^{n-1})\cdot(h_0 + h_1 X + \dots + X^k),
\end{align*}
Beim Nullpolynom verschwinden sämtliche Koeffizienten, insbesondere von $X^k$, $X^{k+1}$, \dots, $X^{n-1}$.

Koeffizientenvergleich für $X^k$
\begin{align*}
0 &= c_0 + c_1 h_{k-1} + \dots + c_k h_0 \\
&= (1\ h_{k-1} \dots h_0)\,(c_0\ c_1 \dots c_k)^T \\
&= (1\ h_{k-1} \dots h_0\ 0 \dots 0)\,(c_0\ c_1 \dots c_{n-1})^T
\end{align*}
Koeffizientenvergleich für $X^{k+1}$
\begin{align*}
0 &= c_1 + c_2 h_{k-1} + \dots + c_{k+1} h_0 \\
&= (1\ h_{k-1} \dots h_0)\,(c_1\ c_2 \dots c_{k+1})^T \\
&= (0\ 1\ h_{k-1} \dots h_0\ 0 \dots 0)\,(c_0\ c_1 \dots c_{n-1})^T
\end{align*}
\dots

Koeffizientenvergleich für $X^{n-1}$
\begin{align*}
0 &= c_{r-1} + c_r h_{k-1} + \dots + c_{n-1} h_0 \\
&= (1\ h_{k-1} \dots h_0)\,(c_{r-1}\ c_r \dots c_{n-1})^T \\
&= (0 \dots 0\ 1\ h_{k-1} \dots h_0)\,(c_0\ c_1 \dots c_{n-1})^T
\end{align*}

Matrixschreibweise
\[
\begin{pmatrix}
1 & h_{k-1} & \cdots & h_1 & h_0 & 0 & \cdots & 0 \\
0 & 1 & h_{k-1} & \cdots & h_1 & h_0 & & 0 \\
 & \ddots & \ddots & \ddots & & \ddots & \ddots & \\
0 & & 0 & 1 & h_{k-1} & \cdots & h_1 & h_0
\end{pmatrix}
\begin{pmatrix}
c_0 \\ c_1 \\ \vdots \\ c_{n-1}
\end{pmatrix}
= H\, c^T \equiv \mathbf{o}
\]

\begin{definition}[reziprokes Polynom]
\[
\overleftarrow{h}(X) := 1 + h_{k-1} X + \dots + h_0 X^k = X^k\, h(X^{-1})
\]
\end{definition}

Kontrollmatrix
\[
H = \begin{pmatrix}
\overleftarrow{h}(X) \\ X * \overleftarrow{h}(X) \\ \vdots \\ X^{r-1} * \overleftarrow{h}(X)
\end{pmatrix}
\]

\begin{beispiel}
$C_2$: \quad $G = H = \begin{pmatrix} 1 & 0 & 1 & 0 \\ 0 & 1 & 0 & 1 \end{pmatrix}$
\end{beispiel}

\begin{korollar}[Folgerung]
Ein zyklischer Code $C$ besitzt eine zyklische Generator- und Kontrollmatrix.
\end{korollar}

\subsection{5.6* Der binäre Golay Code $G_{23}$}

Zerlegen $X^{23} - 1 \in F_2[X]$ in irreduzible Polynome.\\
Die Kreisteilungsklassen sind
\begin{align*}
\Zm{23} = \{0\} &\cup \{1, 2, 4, 8, -7, 9, -5, -10, 3, 6, 12\} \\
&\cup \{-1, -2, -4, -8, 7, -9, 5, 10, -3, -6, -12\}
\end{align*}
sagt uns, dass $X^{23} - 1$ in drei irreduzible (unzerlegbare) Polynome zerfällt
\[
\Rightarrow\quad X^{23} - 1 = (X+1)\, g_1(X)\, g_2(X), \quad \operatorname{grad} g_1(X) = \operatorname{grad} g_2(X) = 11
\]

Koeffizientenvergleich und Lösen des Gleichungssystems liefert
\begin{align*}
g_1(X) &= 1 + X^2 + X^4 + X^5 + X^6 + X^{10} + X^{11} \\
g_2(X) &= 1 + X + X^5 + X^6 + X^7 + X^9 + X^{11} = X^{11}\, g_1(X^{-1}) = \overleftarrow{g_1}(X)
\end{align*}

\begin{bemerkung}[Hinweis]
Eine Zerlegung $X^n - 1 = (X+1)\, g_1(X)\, \overleftarrow{g_1}(X)$ gibt es, wenn $n$ Primzahl und die Form $8m - 1$ besitzt.
\end{bemerkung}

\begin{definition}
Der binäre zyklische Code mit Wortlänge $n = 23$ und Generatorpolynom
\[
g(X) = 1 + X^2 + X^4 + X^5 + X^6 + X^{10} + X^{11}
\]
heißt binärer \textbf{Golay Code $G_{23}$}.
\end{definition}

\begin{tabular}{ll}
Dim $G_{23}$ & $k = n - \operatorname{grad} g(X) = 12$ \\
Anzahl der Codewörter & $|G_{23}| = 2^{12} = 4096$
\end{tabular}

Die Erzeugung dieser 4096 Codewörter $\mathbf{c} \in G_{23}$ und Ermittlung ihrer Gewichte $w(\mathbf{c})$ liefert das Minimalgewicht
\[
d = 7
\]

\begin{satz}
$G_{23}$ ist ein perfekter Code mit den Parametern $[23, 12, 7]_2$.
\end{satz}

\begin{proof}
Bleibt zu zeigen, dass $G_{23}$ die Hamming-Schranke mit \glqq$=$\grqq{} erfüllt.

\textbf{Hamming-Schranke}\quad $|C| \cdot \displaystyle\sum_{i=0}^{e} \binom{n}{i} (q-1)^i \le q^n$, \quad $2e + 1 = d$
\begin{align*}
|G_{23}| \cdot \sum_{i=0}^{e} \binom{n}{i} (q-1)^i &= 2^{12} \cdot \sum_{i=0}^{3} \binom{23}{i} (2-1)^i \\
&= 2^{12} \cdot \Bigl(1 + 23 + \frac{23\cdot 22}{1\cdot 2} + \frac{23\cdot 22\cdot 21}{1\cdot 2\cdot 3}\Bigr) \\
&= 2^{12} \cdot (1 + 23 + 253 + 1771) \\
&= 2^{12} \cdot 2048 = 2^{12} \cdot 2^{11} = 2^{23} \qedhere
\end{align*}
\end{proof}

\subsection{5.7* Der ternäre Golay Code $G_{11}$}

Zerlegen $X^{11} - 1 \in F_3[X]$ in irreduzible Polynome.\\
Die Kreisteilungsklassen sind
\[
\Zm{11} = \{0\} \cup \{1, 3, -2, 5, 4\} \cup \{-1, -3, 2, -5, -4\}
\]
\begin{align*}
\Rightarrow\quad & X^{11} - 1 = (X - 1)\, g_1(X)\, g_2(X), \quad \operatorname{grad} g_1(X) = \operatorname{grad} g_2(X) = 5 \\
\Rightarrow\quad & X^{10} + X^9 + \dots + X + 1 = g_1(X)\, g_2(X)
\end{align*}

Koeffizientenvergleich und Lösen des Gleichungssystems liefert
\begin{align*}
g_1(X) &= -1 + X^2 - X^3 + X^4 + X^5 \\
g_2(X) &= -1 - X + X^2 - X^3 + X^5 = -X^5\, g_1(X^{-1}) = -\overleftarrow{g_1}(X)
\end{align*}
z.B.\ $X^2$: \qquad $1 = -1 - 1 = -2 = 1$ \quad in $\mathbf{F}_3$

\begin{definition}
Der ternäre zyklische Code mit Wortlänge $n = 11$ und Generatorpolynom
\[
g(X) = -1 + X^2 - X^3 + X^4 + X^5
\]
heißt ternärer \textbf{Golay Code $G_{11}$}.
\end{definition}

\begin{tabular}{ll}
Dimension von $G_{11}$ & $k = \operatorname{Dim} G_{11} = n - \operatorname{grad} g(X) = 6$ \\
Anzahl der Codewörter & $|G_{11}| = 3^6 = 729$
\end{tabular}

Die Erzeugung dieser 729 Codewörter $\mathbf{c} \in G_{11}$ und Ermittlung ihrer Gewichte $w(\mathbf{c})$ liefert das Minimalgewicht
\[
d = 5
\]

\begin{satz}
$G_{11}$ ist ein perfekter Code mit den Parametern $[11, 6, 5]_3$.
\end{satz}

\begin{proof}
Bleibt zu zeigen, dass $G_{11}$ die Hamming-Schranke mit \glqq$=$\grqq{} erfüllt.

\textbf{Hamming-Schranke}\quad $|C| \cdot \displaystyle\sum_{i=0}^{e} \binom{n}{i} (q-1)^i \le q^n$, \quad $2e + 1 = d$
\begin{align*}
|G_{11}| \cdot \sum_{i=0}^{e} \binom{n}{i} (q-1)^i &= 3^6 \cdot \sum_{i=0}^{2} \binom{11}{i} (3-1)^i \\
&= 3^6 \cdot \Bigl(1 + 11 \cdot 2 + \frac{11\cdot 10}{1\cdot 2} \cdot 2^2\Bigr) \\
&= 3^6 \cdot (1 + 22 + 55 \cdot 4) = 3^6 \cdot 243 = 3^{11} \qedhere
\end{align*}
\end{proof}
